\EMsection{روش دالامبر در حل معادلهٔ موجی یک‌بعدی}{روش دالامبر در حل معادلهٔ موجی یک‌بعدی}

\EMdisp{\begin{cases}\dfrac{\partial^2u}{\partial t^2}=c^2\dfrac{\partial^2u}{\partial x^2}\\\noalign{\vskip 2mm} u(0,t)=u(L,t)=0\\ u(x,0)=f(x)\\ u_t(x,0)=g(x)\end{cases}}
با استفاده از تغییر متغیرها به شکل زیر معادلهٔ دیفرانسیل را بازنویسی و حل می‌کنیم.
\EMdisp{u(x,t)\EMbr \Rightarrow u(\varphi,\psi),\qquad\EMbr \begin{cases}\varphi=x+ct\\ \psi=x-ct\end{cases}}
\EMdisp{\frac{\partial u}{\partial x}\EMbr =\frac{\partial u}{\partial\varphi}\frac{\partial\varphi}{\partial x}+\frac{\partial u}{\partial\psi}\frac{\partial\psi}{\partial x}\EMbr =\frac{\partial u}{\partial\varphi}+\frac{\partial u}{\partial\psi}}
\EMdisp{\frac{\partial u}{\partial t}\EMbr =\frac{\partial u}{\partial\varphi}\frac{\partial\varphi}{\partial t}+\frac{\partial u}{\partial\psi}\frac{\partial\psi}{\partial t}\EMbr =c\frac{\partial u}{\partial\varphi}-c\frac{\partial u}{\partial\psi}}
\EMdisp{\frac{\partial^2u}{\partial x^2}\EMbr =\frac\partial{\partial x}\left(\frac{\partial u}{\partial\varphi}+\frac{\partial u}{\partial\psi}\right)\EMbr =\left(\frac{\partial^2u}{\partial\varphi^2}+\frac{\partial^2u}{\partial\varphi\partial\psi}\right)\cdot1+\left(\frac{\partial^2u}{\partial\psi\partial\varphi}+\frac{\partial^2u}{\partial\psi^2}\right)\cdot1\EMbr =\frac{\partial^2u}{\partial\varphi^2}+2\frac{\partial^2u}{\partial\varphi\partial\psi}+\frac{\partial^2u}{\partial\psi^2}}
(با استفاده از \EMim{\dfrac{\partial^2u}{\partial\varphi\partial\psi}=\dfrac{\partial^2u}{\partial\psi\partial\varphi}}.)

\begin{EMexercise}{}

به طریق مشابه تساوی زیر را ثابت کنید:
\EMdisp{\frac{\partial^2u}{\partial t^2}\EMbr =c^2\left(\frac{\partial^2u}{\partial\varphi^2}-2\frac{\partial^2u}{\partial\varphi\partial\psi}+\frac{\partial^2u}{\partial\psi^2}\right)}

\end{EMexercise}

با قرار دادن دو رابطهٔ بالا در معادلهٔ دیفرانسیل موجی یک‌بعدی:
\EMdisp{c^2\left(\frac{\partial^2u}{\partial\varphi^2}-2\frac{\partial^2u}{\partial\varphi\partial\psi}+\frac{\partial^2u}{\partial\psi^2}\right)\EMbr =c^2\left(\frac{\partial^2u}{\partial\varphi^2}+2\frac{\partial^2u}{\partial\varphi\partial\psi}+\frac{\partial^2u}{\partial\psi^2}\right)\quad\EMbr \EMbr \Longrightarrow\quad\EMbr \frac{\partial^2u}{\partial\varphi\partial\psi}\EMbr =0}

\textbf{حل عمومی:}
\EMdisp{\frac{\partial u}{\partial\varphi}\EMbr =\phi_1(\varphi)\ \EMbr \Rightarrow\ u(\varphi,\psi)\EMbr =\int\phi_1(\varphi)\,d\varphi+\Psi(\psi)\EMbr =\Phi(\varphi)+\Psi(\psi)\quad\EMbr \EMbr \Longrightarrow\quad\EMbr  u(x,t)\EMbr =\Phi(x+ct)+\Psi(x-ct)}

\textbf{اعمال شرایط اولیه:}
\EMdisp{u(x,0)\EMbr =\Phi(x)+\Psi(x)\EMbr =f(x)}
\EMdisp{u_t(x,t)\EMbr =\frac{\partial\Phi}{\partial\varphi}\frac{\partial\varphi}{\partial t}+\frac{\partial\Psi}{\partial\psi}\frac{\partial\psi}{\partial t}\EMbr =c\Phi'(\varphi)-c\Psi'(\psi)\EMbr =c\Phi'(x+ct)-c\Psi'(x-ct)}
\EMdisp{u_t(x,0)\EMbr =c\Phi'(x)-c\Psi'(x)\EMbr =c\big[\Phi'(x)-\Psi'(x)\big]\EMbr =g(x)}
\EMdisp{\begin{cases}\Phi(x)+\Psi(x)=f(x)\\ \Phi'(x)-\Psi'(x)=\dfrac1cg(x)\end{cases}\EMbr \Rightarrow\begin{cases}\Phi'(x)+\Psi'(x)=f'(x)\\ \Phi'(x)-\Psi'(x)=\dfrac1cg(x)\end{cases}\EMbr \Rightarrow\begin{cases}\Phi'(x)=\dfrac12f'(x)+\dfrac1{2c}g(x)\\\noalign{\vskip 2mm} \Psi'(x)=\dfrac12f'(x)-\dfrac1{2c}g(x)\end{cases}}
\EMdisp{\begin{cases}\Phi(x)=\dfrac12f(x)+\dfrac1{2c}\displaystyle\int_0^xg(s)\,ds+k_1\\\noalign{\vskip 2mm} \Psi(x)=\dfrac12f(x)-\dfrac1{2c}\displaystyle\int_0^xg(s)\,ds+k_2\end{cases}}
\EMdisp{u(x,t)\EMbr =\frac12\left(f(x+ct)+\frac1c\int_0^{x+ct}g(s)\,ds\right)+\frac12\left(f(x-ct)-\frac1c\int_0^{x-ct}g(s)\,ds\right)+k_1+k_2}

\begin{EMimportant}{فرمول دالامبر}

\EMdisp{u(x,t)\EMbr =\frac12\Big[f(x+ct)+f(x-ct)\Big]+\frac1{2c}\int_{x-ct}^{x+ct}g(s)\,ds+K,\qquad\EMbr  u(0,0)\EMbr =0\ \EMbr \Rightarrow\ K\EMbr =0}
\EMdisp{u(x,t)\EMbr =\frac12\Big[f(x+ct)+f(x-ct)\Big]+\frac1{2c}\int_{x-ct}^{x+ct}g(s)\,ds}

\end{EMimportant}

\begin{EMexercise}{}

ثابت کنید که اگر بخواهیم شرایط مرزی برقرار باشند بایستی در رابطهٔ بالا توابع \EMim{f(x)} و \EMim{g(x)} هر دو تناوبی با دورهٔ تناوب \EMim{2L} بوده و فرد باشند.

\end{EMexercise}

\EMsection{معادلهٔ دیفرانسیل انتقال حرارت}{معادلهٔ دیفرانسیل انتقال حرارت}

\EMdisp{\frac{\partial u}{\partial t}\EMbr =c^2\frac{\partial^2u}{\partial x^2},\qquad\EMbr  \frac{\partial u}{\partial t}\EMbr =c^2\left(\frac{\partial^2u}{\partial x^2}+\frac{\partial^2u}{\partial y^2}\right),\qquad\EMbr  \frac{\partial u}{\partial t}\EMbr =c^2\left(\frac{\partial^2u}{\partial x^2}+\frac{\partial^2u}{\partial y^2}+\frac{\partial^2u}{\partial z^2}\right)}

\EMkeepfig{m12-rod}{3}

\EMsubsection{معادلهٔ دیفرانسیل انتقال حرارت یک‌بعدی}{معادلهٔ دیفرانسیل انتقال حرارت یک‌بعدی}

\EMfigure{m12-rod}{میلهٔ یک‌بعدی \EMim{0\le x\le L} با توزیع دمای \EMim{u(x,t)}}

برای حل یکتای این معادله نیاز به شرایط مرزی و شرایط اولیه داریم:

\begin{itemize}
\tightlist
\item
  شرایط مرزی (\lr{Boundary Conditions}): \EMim{u(0,t)=u(L,t)=0}
\item
  شرایط اولیه (\lr{Initial Conditions}): \EMim{u(x,0)=f(x)}
\end{itemize}

\EMsection{حل معادلهٔ دیفرانسیل انتقال حرارت یک‌بعدی}{حل معادلهٔ دیفرانسیل انتقال حرارت یک‌بعدی}

\EMdisp{\begin{cases}\dfrac{\partial u}{\partial t}=c^2\dfrac{\partial^2u}{\partial x^2}\\\noalign{\vskip 2mm} u(0,t)=u(L,t)=0\\ u(x,0)=f(x)\end{cases}}
در اینجا برای حل این معادله از روش \textbf{جداسازی متغیرها} استفاده می‌کنیم. فرض: \EMim{u(x,t)=X(x)\times T(t)}.
\EMdisp{\frac{\partial^2u}{\partial x^2}\EMbr =X''(x)T(t),\qquad\EMbr  \frac{\partial u}{\partial t}\EMbr =X(x)T'(t)\quad\EMbr \EMbr \Longrightarrow\quad\EMbr  X(x)T'(t)\EMbr =c^2X''(x)T(t)\quad\EMbr \EMbr \Longrightarrow\quad\EMbr \frac{T'}{c^2T}\EMbr =\frac{X''}X\EMbr =k}
\EMim{k} باید حتماً یک عدد منفی باشد (برای \EMim{k=0} و \EMim{k>0} دقیقاً مانند مسئلهٔ موجی \EMim{X(x)=0} به دست می‌آید). با \EMim{k=-\lambda^2}:
\EMdisp{\frac{T'}{c^2T}\EMbr =\frac{X''}X\EMbr =-\lambda^2\quad\EMbr \EMbr \Longrightarrow\quad\EMbr \begin{cases}X''+\lambda^2X=0\\ T'+c^2\lambda^2T=0\end{cases}}
\EMdisp{X(x)\EMbr =A\cos\lambda x+B\sin\lambda x,\quad\EMbr  X(0)\EMbr =0\EMbr \Rightarrow A\EMbr =0,\qquad\EMbr  X(L)\EMbr =0\EMbr \Rightarrow B\sin\lambda L\EMbr =0\EMbr \Rightarrow\lambda\EMbr =\frac{n\pi}L\EMbr =\lambda_n}
\EMdisp{X_n(x)\EMbr =B_n\sin\lambda_nx,\qquad\EMbr  T'+c^2\lambda_n^2T\EMbr =0\ \EMbr \Rightarrow\ T_n(t)\EMbr =C_ne^{-(c\lambda_n)^2t}}
\EMdisp{u_n(x,t)\EMbr =X_n(x)\times T_n(t)\EMbr =B_n\sin\!\left(\frac{n\pi}Lx\right)\times C_ne^{-\left(\frac{cn\pi}L\right)^2t}}
\EMdisp{u(x,t)\EMbr =\sum_{n=1}^{\infty}A_ne^{-(c\lambda_n)^2t}\sin(\lambda_nx),\qquad\EMbr  \lambda_n\EMbr =\frac{n\pi}L}
شرط اولیه \EMim{u(x,0)=f(x)}:
\EMdisp{u(x,0)\EMbr =\sum_{n=1}^{\infty}A_n\sin\lambda_nx\EMbr =f(x)}
تابع \EMim{f(x)} با یک تابع تناوبی و فرد برابر شده است، بنابراین برای محاسبهٔ ضرایب مجهول فوریه باید \EMim{f(x)} را یک تابع تناوبی (با دورهٔ تناوب \EMim{2L}) و فرد فرض کنیم:
\EMdisp{A_n\EMbr =\frac2{2L}\int_{-L}^{L}f(x)\sin\!\left(\frac{2\pi}{2L}nx\right)dx\EMbr =\frac1L\int_{-L}^{L}f(x)\sin\!\left(\frac{n\pi}Lx\right)dx\EMbr =\frac2L\int_0^Lf(x)\sin\!\left(\frac{n\pi}Lx\right)dx}

\begin{EMimportant}{جواب معادلهٔ گرمای یک‌بعدی (روش جداسازی متغیرها)}

\EMdisp{u(x,t)\EMbr =\sum_{n=1}^{\infty}A_ne^{-(c\lambda_n)^2t}\sin(\lambda_nx),\qquad\EMbr  \lambda_n\EMbr =\frac{n\pi}L,\qquad\EMbr  A_n\EMbr =\frac2L\int_0^Lf(x)\sin\!\left(\frac{n\pi}Lx\right)dx}

\end{EMimportant}

\begin{EMexample}{مثال ۱}

جواب معادلهٔ انتقال حرارت یک‌بعدی زیر را به‌دست آورید.
\EMdisp{\frac{\partial u}{\partial t}\EMbr =\frac{\partial^2u}{\partial x^2},\qquad\EMbr  u(0,t)\EMbr =u(10,t)\EMbr =0,\qquad\EMbr  u(x,0)\EMbr =f(x)\EMbr =\begin{cases}x&0\le x<5\\ 10-x&5\le x\le10\end{cases}}

\EMdisp{\begin{aligned}
A_n&=\frac2L\int_0^Lf(x)\sin\!\left(\frac{n\pi}Lx\right)dx=0.2\int_0^5x\sin\!\left(\frac{n\pi}{10}x\right)dx+0.2\int_5^{10}(10-x)\sin\!\left(\frac{n\pi}{10}x\right)dx\\
&=0.2\left[(x)\left(\frac{-10}{n\pi}\cos\frac{n\pi}{10}x\right)-(1)\left(\frac{-100}{n^2\pi^2}\sin\frac{n\pi}{10}x\right)\right]_0^5\\
&\quad+0.2\left[(10-x)\left(\frac{-10}{n\pi}\cos\frac{n\pi}{10}x\right)-(-1)\left(\frac{-100}{n^2\pi^2}\sin\frac{n\pi}{10}x\right)\right]_5^{10}=\frac{40}{n^2\pi^2}\sin\frac{n\pi}2
\end{aligned}}
\EMdisp{u(x,t)\EMbr =\sum_{n=1}^{\infty}\left(\frac{40}{n^2\pi^2}\sin\frac{n\pi}2\right)e^{-\left(\frac{cn\pi}{10}\right)^2t}\sin\!\left(\frac{n\pi}{10}x\right)}
(با \EMim{c=1}.)

\end{EMexample}

\begin{EMexample}{مثال ۲ (شرایط مرزی غیرصفر)}

جواب معادلهٔ انتقال حرارت یک‌بعدی زیر را به‌دست آورید.
\EMdisp{\frac{\partial u}{\partial t}\EMbr =c^2\frac{\partial^2u}{\partial x^2},\qquad\EMbr  u(0,t)\EMbr =T_0,\qquad\EMbr  u(L,t)\EMbr =T_1,\qquad\EMbr  u(x,0)\EMbr =f(x)}

\EMdisp{u(x,t)\EMbr =w(x,t)+v(x)\quad\EMbr \EMbr \Longrightarrow\quad\EMbr \begin{cases}w(0,t)+v(0)=T_0\\ w(L,t)+v(L)=T_1\end{cases}}
\EMdisp{\frac{\partial w}{\partial t}\EMbr =c^2\left(\frac{\partial^2w}{\partial x^2}+\frac{\partial^2v}{\partial x^2}\right)}
\EMdisp{\begin{cases}\dfrac{\partial^2v}{\partial x^2}=0\\ v(0)=T_0\\ v(L)=T_1\end{cases}\EMbr \Rightarrow\ v(x)\EMbr =\left(\frac{T_1-T_0}L\right)x+T_0,\qquad\EMbr \begin{cases}\dfrac{\partial w}{\partial t}=c^2\dfrac{\partial^2w}{\partial x^2}\\ w(0,t)=w(L,t)=0\end{cases}\EMbr \Rightarrow\ w(x,t)\EMbr =\sum_{n=1}^{\infty}A_ne^{-(c\lambda_n)^2t}\sin(\lambda_nx)}
\EMdisp{u(x,t)\EMbr =\underbrace{\sum_{n=1}^{\infty}A_ne^{-(c\lambda_n)^2t}\sin(\lambda_nx)}_{w(x,t)}+\underbrace{\left(\frac{T_1-T_0}L\right)x+T_0}_{v(x)},\qquad\EMbr  \lambda_n\EMbr =\frac{n\pi}L}
شرط اولیه:
\EMdisp{\sum_{n=1}^{\infty}A_n\sin\frac{n\pi}Lx\EMbr =f(x)-\left(\frac{T_1-T_0}L\right)x-T_0\triangleq f_1(x)\ \EMbr \Longrightarrow\ A_n\EMbr =\frac2L\int_0^L\left[f(x)-\left(\frac{T_1-T_0}L\right)x-T_0\right]\sin\frac{n\pi}Lx\,dx}

\end{EMexample}

\begin{EMexample}{مثال ۳ (مرزهای عایق)}

جواب معادلهٔ انتقال حرارت یک‌بعدی زیر را به‌دست آورید.
\EMdisp{\frac{\partial u}{\partial t}\EMbr =c^2\frac{\partial^2u}{\partial x^2},\qquad\EMbr  u_x(0,t)\EMbr =u_x(L,t)\EMbr =0,\qquad\EMbr  u(x,0)\EMbr =f(x)}

برای حل این مسئله از روش \textbf{جداسازی متغیرها} استفاده می‌کنیم: \EMim{u(x,t)=X(x)\times T(t)}.
\EMdisp{\frac{T'}{c^2T}\EMbr =\frac{X''}X\EMbr =-\lambda^2\ \EMbr \Rightarrow\ \begin{cases}X''+\lambda^2X=0\ \Rightarrow\ X(x)=A\cos\lambda x+B\sin\lambda x\\ T'+c^2\lambda^2T=0\end{cases}}
\EMdisp{X'(x)\EMbr =-A\lambda\sin\lambda x+B\lambda\cos\lambda x,\qquad\EMbr  X'(0)\EMbr =0\EMbr \Rightarrow B\EMbr =0,\qquad\EMbr  X'(L)\EMbr =0\EMbr \Rightarrow-A\lambda\sin\lambda L\EMbr =0\EMbr \Rightarrow\lambda L\EMbr =n\pi\EMbr \Rightarrow\lambda_n\EMbr =\frac{n\pi}L}
\EMdisp{X_n(x)\EMbr =A_n\cos\lambda_nx,\qquad\EMbr  T_n(t)\EMbr =C_ne^{-(c\lambda_n)^2t}}
\EMdisp{u(x,t)\EMbr =\sum_{n=1}^{\infty}B_ne^{-(c\lambda_n)^2t}\cos(\lambda_nx),\qquad\EMbr  \lambda_n\EMbr =\frac{n\pi}L}
شرط اولیه \EMim{u(x,0)=\sum_{n=1}^{\infty}B_n\cos\!\left(\dfrac{n\pi}Lx\right)=f(x)}:
\EMdisp{B_n\EMbr =\frac2L\int_0^Lf(x)\cos\!\left(\frac{n\pi}Lx\right)dx}

\end{EMexample}

\EMsection{معادلهٔ انتقال حرارت یک‌بعدی بر روی مفتول نامتناهی}{معادلهٔ انتقال حرارت یک‌بعدی بر روی مفتول نامتناهی}

\EMdisp{\begin{cases}\dfrac{\partial u}{\partial t}=c^2\dfrac{\partial^2u}{\partial x^2}\\\noalign{\vskip 2mm} u(\infty,t)=u(-\infty,t)=0\\ u(x,0)=f(x),\quad-\infty<x<\infty\end{cases}}
با فرض \EMim{u(x,t)=X(x)\times T(t)}:
\EMdisp{X(x)T'(t)\EMbr =c^2X''(x)T(t)\ \EMbr \Rightarrow\ \frac{T'}{c^2T}\EMbr =\frac{X''}X\EMbr =-\lambda^2\ \EMbr \Rightarrow\ \begin{cases}X(x)=A_1(\lambda)\cos\lambda x+B_1(\lambda)\sin\lambda x=X_\lambda(x)\\ T(t)=C_1(\lambda)e^{-c^2\lambda^2t}=T_\lambda(t)\end{cases}}
\EMdisp{u_\lambda(x,t)\EMbr =\big[A(\lambda)\cos\lambda x+B(\lambda)\sin\lambda x\big]e^{-c^2\lambda^2t}\quad\EMbr \EMbr \Longrightarrow\quad\EMbr  u(x,t)\EMbr =\frac1\pi\int_0^\infty\big[A(\lambda)\cos\lambda x+B(\lambda)\sin\lambda x\big]e^{-c^2\lambda^2t}\,d\lambda}
صحت معادله:
\EMdisp{\frac{\partial^2u}{\partial x^2}\EMbr =\frac1\pi\int_0^\infty-\lambda^2\big[A\cos\lambda x+B\sin\lambda x\big]e^{-c^2\lambda^2t}\,d\lambda,\qquad\EMbr  \frac{\partial u}{\partial t}\EMbr =\frac1\pi\int_0^\infty-c^2\lambda^2\big[A\cos\lambda x+B\sin\lambda x\big]e^{-c^2\lambda^2t}\,d\lambda\ \EMbr \Rightarrow\ \frac{\partial u}{\partial t}\EMbr =c^2\frac{\partial^2u}{\partial x^2}\ \EMcheck }
\textbf{اعمال شرط اولیه:} \EMim{u(x,0)=\dfrac1\pi\displaystyle\int_0^\infty\big[A(\lambda)\cos\lambda x+B(\lambda)\sin\lambda x\big]d\lambda=f(x)}، یعنی انتگرال فوریه (فصل ۱۰):
\EMdisp{A(\lambda)\EMbr =\mathfrak{I}_C\{f(x)\}\EMbr =\int_{-\infty}^{\infty}f(x)\cos\lambda x\,dx,\qquad\EMbr  B(\lambda)\EMbr =\mathfrak{I}_S\{f(x)\}\EMbr =\int_{-\infty}^{\infty}f(x)\sin\lambda x\,dx}
\EMdisp{\begin{aligned}
u(x,t)&=\frac1\pi\int_0^\infty\left[\left(\int_{-\infty}^{\infty}f(\nu)\cos\lambda\nu\,d\nu\right)\cos\lambda x+\left(\int_{-\infty}^{\infty}f(\nu)\sin\lambda\nu\,d\nu\right)\sin\lambda x\right]e^{-c^2\lambda^2t}\,d\lambda\\
&=\frac1\pi\int_0^\infty\left\{\int_{-\infty}^{\infty}f(\nu)\big[\cos\lambda\nu\cos\lambda x+\sin\lambda\nu\sin\lambda x\big]d\nu\right\}e^{-c^2\lambda^2t}\,d\lambda\\
&=\frac1\pi\int_0^\infty\left\{\int_{-\infty}^{\infty}f(\nu)\cos\lambda(x-\nu)\,d\nu\right\}e^{-c^2\lambda^2t}\,d\lambda
\end{aligned}}

\begin{EMimportant}{شکل اول جواب}

\EMdisp{u(x,t)\EMbr =\frac1\pi\int_{-\infty}^{\infty}f(\nu)\left\{\int_0^\infty e^{-c^2\lambda^2t}\cos\lambda(x-\nu)\,d\lambda\right\}d\nu}

\end{EMimportant}

\EMhold

\EMsubsection{لم‌های ریاضی}{لم‌های ریاضی}

\begin{EMtheorem}{لم (الف)}

\EMdisp{I\EMbr =\int_0^\infty e^{-x^2}\,dx\EMbr =\frac{\sqrt\pi}2}

\end{EMtheorem}

\begin{EMproof}

\EMdisp{I^2\EMbr =\int_0^\infty e^{-x^2}dx\times\int_0^\infty e^{-y^2}dy\EMbr =\int_0^\infty\!\!\int_0^\infty e^{-(x^2+y^2)}\,dx\,dy\ \overset{\rho^2=x^2+y^2}{=}\ \int_0^{\pi/2}\!\!\int_0^\infty\rho\,e^{-\rho^2}\,d\rho\,d\theta\EMbr =\int_0^{\pi/2}\left[\frac{-1}2e^{-\rho^2}\right]_0^\infty d\theta\EMbr =\frac\pi4}

\end{EMproof}

\begin{EMtheorem}{لم (ب)}

\EMdisp{\int_0^\infty e^{-s^2}\cos(2bs)\,ds\EMbr =\frac{\sqrt\pi}2e^{-b^2}}

\end{EMtheorem}

\begin{EMproof}

\EMdisp{I(b)\EMbr =\int_0^\infty e^{-s^2}\cos(2bs)\,ds\ \EMbr \Rightarrow\ \frac{dI(b)}{db}\EMbr =-\int_0^\infty2se^{-s^2}\sin(2bs)\,ds}
با انتگرال‌گیری جزء به جزء (\EMim{du=-2se^{-s^2}ds}، \EMim{v=\sin2bs}):
\EMdisp{\frac{dI(b)}{db}\EMbr =e^{-s^2}\sin2bs\Big|_0^\infty-2b\int_0^\infty e^{-s^2}\cos2bs\,ds\EMbr =-2b\,I(b)\ \EMbr \Rightarrow\ I(b)\EMbr =ke^{-b^2},\qquad\EMbr  I(0)\EMbr =k\EMbr =\int_0^\infty e^{-s^2}ds\EMbr =\frac{\sqrt\pi}2}

\end{EMproof}

با تغییر متغیر \EMim{c\lambda\sqrt t=s}:
\EMdisp{\int_0^\infty e^{-c^2\lambda^2t}\cos\lambda(x-\nu)\,d\lambda\EMbr =\frac1{c\sqrt t}\int_0^\infty e^{-s^2}\cos\!\left[2\frac{(x-\nu)}{2c\sqrt t}s\right]ds\EMbr =\frac1{c\sqrt t}\frac{\sqrt\pi}2e^{-\left(\frac{x-\nu}{2c\sqrt t}\right)^2}}

\begin{EMimportant}{شکل دوم جواب}

\EMdisp{u(x,t)\EMbr =\frac1{2c\sqrt{\pi t}}\int_{-\infty}^{\infty}f(\nu)\,e^{-\frac{(x-\nu)^2}{4c^2t}}\,d\nu}

\end{EMimportant}

با تغییر متغیر \EMim{z=\dfrac{\nu-x}{2c\sqrt t}} (یعنی \EMim{\nu=x+2c\sqrt t\,z} و \EMim{d\nu=2c\sqrt t\,dz}):

\begin{EMimportant}{شکل سوم جواب}

\EMdisp{u(x,t)\EMbr =\frac1{\sqrt\pi}\int_{-\infty}^{\infty}f\big(x+2c\sqrt t\,z\big)\,e^{-z^2}\,dz}

\end{EMimportant}

\EMsection{تابع خطا}{تابع خطا}

در ریاضی اغلب تابع خطا را به صورت زیر تعریف می‌کنند:
\EMdisp{\operatorname{erf}(x)\triangleq\frac2{\sqrt\pi}\int_0^xe^{-t^2}\,dt}
خواص تابع خطا:

\begin{itemize}
\tightlist
\item
  خاصیت تقارن فرد: \EMim{\operatorname{erf}(-x)=-\operatorname{erf}(x)}
\item
  خاصیت مجانبی: \EMim{\operatorname{erf}(\infty)=1}
\item
  بسط مک‌لورن این تابع به صورت زیر است (با بسط \EMim{e^x} شروع و این تساوی را ثابت کنید):
  \EMdisp{\operatorname{erf}(x)\EMbr =\frac2{\sqrt\pi}\left[x-\frac{x^3}{1!\times3}+\frac{x^5}{2!\times5}-\frac{x^7}{3!\times7}+-\cdots\right]\EMbr =\frac2{\sqrt\pi}\sum_{n=1}^{\infty}(-1)^{n+1}\frac{x^{2n-1}}{(n-1)!\,(2n-1)}}
\end{itemize}

\EMfigure{m12-erf}{نمودار تابع خطا \EMim{\operatorname{erf}(x)} برای \EMim{-2\le x\le2}}

\begin{EMexample}{مثال ۱}

جواب معادلهٔ بالا را با فرض شرط اولیهٔ مقابل به دست آورید.
\EMdisp{f(x)\EMbr =\begin{cases}U_0&|x|<1\\ 0&|x|\ge1\end{cases}}

از شکل دوم جواب استفاده می‌کنیم:
\EMdisp{u(x,t)\EMbr =\frac{U_0}{2c\sqrt{\pi t}}\int_{-1}^{1}e^{-\frac{(x-\nu)^2}{4c^2t}}\,d\nu\ \overset{z=\frac{\nu-x}{2c\sqrt t}}{=}\ \frac{U_0}{\sqrt\pi}\int_{-\frac{1+x}{2c\sqrt t}}^{\frac{1-x}{2c\sqrt t}}e^{-z^2}\,dz}
\EMdisp{=\frac{U_0}{\sqrt\pi}\left\{\int_{-\frac{1+x}{2c\sqrt t}}^{0}e^{-z^2}dz+\int_0^{\frac{1-x}{2c\sqrt t}}e^{-z^2}dz\right\}\EMbr =\frac{U_0}{\sqrt\pi}\frac{\sqrt\pi}2\left\{\frac2{\sqrt\pi}\int_0^{\frac{1+x}{2c\sqrt t}}e^{-z^2}dz+\frac2{\sqrt\pi}\int_0^{\frac{1-x}{2c\sqrt t}}e^{-z^2}dz\right\}}
\EMdisp{u(x,t)\EMbr =\frac{U_0}2\left[\operatorname{erf}\!\left(\frac{1+x}{2c\sqrt t}\right)+\operatorname{erf}\!\left(\frac{1-x}{2c\sqrt t}\right)\right]}

\end{EMexample}

\EMfigure{m12-heat-pulse}{توزیع دما \EMim{u(x,t)} در لحظه‌های \EMim{t=0}، \EMim{\frac18}،
\EMim{\frac12}، \EMim{1}، \EMim{2} و \EMim{8} برای پالس اولیهٔ
\EMim{U_0=100} و \EMim{c=1}}

\begin{EMexample}{مثال ۲}

جواب معادلهٔ بالا را با فرض شرط اولیهٔ مقابل به دست آورید.
\EMdisp{f(x)\EMbr =\begin{cases}1&x>0\\ 0&x\le0\end{cases}}

از شکل دوم جواب استفاده می‌کنیم:
\EMdisp{u(x,t)\EMbr =\frac1{2c\sqrt{\pi t}}\int_0^\infty e^{-\frac{(x-\nu)^2}{4c^2t}}\,d\nu\ \overset{z=\frac{\nu-x}{2c\sqrt t}}{=}\ \frac1{\sqrt\pi}\int_{-\frac x{2c\sqrt t}}^{\infty}e^{-z^2}\,dz\EMbr =\frac1{\sqrt\pi}\left\{\int_{-\frac x{2c\sqrt t}}^{0}e^{-z^2}dz+\int_0^\infty e^{-z^2}dz\right\}}
\EMdisp{u(x,t)\EMbr =\frac12\left[\operatorname{erf}\!\left(\frac x{2c\sqrt t}\right)+1\right]}

\end{EMexample}
